\documentclass[11pt]{article}

\setlength{\parindent}{0pt}
\usepackage{xltxtra}
\usepackage{hyperref}
\hypersetup{hidelinks}
\usepackage{url}
\urlstyle{tt}
\usepackage{xcolor}
\definecolor{CVBlue}{RGB}{23,110,191}
\usepackage{calc}
\usepackage{graphicx}
\usepackage{tikz}
\usetikzlibrary{calc}
\usepackage{fontspec}
\usepackage{xeCJK}
\usepackage{enumitem}
\CJKsetecglue{} %% 取消中文与数字之间的间隙

%% 主文档字体设置
\setmainfont[
    Path = fonts/Main/,
    Extension = .otf,
    BoldFont = texgyretermes-bold.otf, % 加粗字体
]{texgyretermes-regular.otf} % 正文字体

% 中文字体设置
\setCJKmainfont[
    Path = fonts/hansans/,
    Extension = .ttf,
    BoldFont = NotoSansSC-Bold.ttf, % 加粗字体
]{NotoSansSC-Regular.otf} % 正文字体

\usepackage{relsize}
\usepackage{xspace}

% 使用 fontawesome（部分图标）
\usepackage{fontawesome} 

% A4纸，上下左右边距
\usepackage[
    a4paper,
    left=1.2cm,
    right=1.2cm,
    top=1.5cm,
    bottom=1cm,
    nohead
]{geometry}

\renewcommand{\baselinestretch}{1.5} % 行间距设为1.5

\usepackage{titlesec}
\usepackage{enumitem}
\setlist{noitemsep} % 取消列表项间的额外间距
\setlist[itemize]{topsep=0.25em, leftmargin=*}
\setlist[enumerate]{topsep=0.25em, leftmargin=*}

% --- 用于控制【不同项目之间】的垂直距离 ---
\newlength{\interProjectSpacing}
\setlength{\interProjectSpacing}{0.9em}
\newcommand{\projectsep}{\vspace{\interProjectSpacing}}

% --- 用于控制【项目标题】与下方【项目描述】的距离 ---
\newlength{\intraProjectTitleSep}
\setlength{\intraProjectTitleSep}{0.4em}
\newcommand{\titlebreak}{\\[\intraProjectTitleSep]}

% --- 用于控制【项目描述】与下方【要点列表】的距离 ---
\newlength{\intraProjectListTopSep}
\setlength{\intraProjectListTopSep}{0.2em}

% =======================================================================

\titleformat{\section}
{\large\bfseries\raggedright}
{}{0em}
{}
[{\color{CVBlue}\titlerule}]
\titlespacing*{\section}{0cm}{*1.6}{*1.2}

\begin{document}
\pagenumbering{gobble}

%%%% 头像与Logo
\begin{tikzpicture}[remember picture, overlay] 
    \node[anchor = north east] at ($(current page.north east)+(-2cm,-1.2cm)$) {\includegraphics[height=3cm]{avatar.jpg}};
\end{tikzpicture}%
\begin{tikzpicture}[remember picture, overlay] 
    \node[anchor = north west] at ($(current page.north west)+(0.5cm,+1.0cm)$) {\includegraphics[height=1.5cm]{ecust.png}};
\end{tikzpicture}%

\centerline{\LARGE\bfseries{徐凯赟}} 
\centerline{\normalsize{求职意向：人工智能 / 计算机视觉 / 深度学习算法工程师}}
\centerline{\normalsize{\faPhone\ *** \quad \faEnvelopeO\ \href{mailto:18321775619@163.com}{18321775619@163.com}}} 
\centerline{\normalsize{\faGithubSquare\ \href{https://github.com/guoshenghuo}{https://github.com/guoshenghuo}}} 

\section{\makebox[\widthof{\faGraduationCap}][c]{\color{CVBlue}\faGraduationCap}\ 教育背景}    
\textbf{华东理工大学} \hfill 2023.09 -- 2027.06（预期）\\[0.5em]
信息工程专业 本科
\begin{itemize}[nosep]
    \item GPA：3.76/4.0，专业排名前10%
    \item 核心课程：机器学习、数字图像处理、信号处理
\end{itemize}

\section{\makebox[\widthof{\faUsers}][c]{\color{CVBlue}\faUsers}\ 项目/科研经历}

% 项目1
\textbf{基于图卷积网络与合成数据的艺术风格抄袭判定} \hfill 2025.11 \titlebreak
剑桥大学项目 | PyTorch, GCN, 可解释性分析
\begin{itemize}[nosep, topsep=\intraProjectListTopSep]
    \item 设计语义注意力对比图，实现AI生成图像的视觉相似度可解释性分析
    \item 搭建基于GCN+MLP的图像相似度学习框架，完成pairwise风格判别
    \item 在ImageNet子集上训练评估，使用Stable Diffusion生成多风格画作并输出相似度评分
    \item 负责模型设计、训练、实验验证与代码实现
\end{itemize}

\projectsep

% 项目2
\textbf{视频说话人转录与字幕对齐系统} \hfill 2025.07 -- 2025.08 \titlebreak
UBC暑期科研项目 | Whisper, 人脸检测, 多模态对齐
\begin{itemize}[nosep, topsep=\intraProjectListTopSep]
    \item 搭建端到端视频转文字pipeline，支持多语音场景
    \item 基于Whisper实现语音识别与语种检测
    \item 结合人脸检测与聚类算法，完成字幕-说话人自动对齐
    \item 负责全流程开发、算法集成与效果优化
\end{itemize}

\projectsep

% 项目3
\textbf{基于3D高斯泼溅的三维建模与生成} \hfill 2025.09 -- 2026.05 \titlebreak
线上科研项目 | 3D Gaussian Splatting, FLAME, 多视角扩散
\begin{itemize}[nosep, topsep=\intraProjectListTopSep]
    \item 基于多视角扩散模型先验，构建高保真人脸重建与虚拟人生成系统
    \item 降低单张图片重建细节损失，融合FLAME参数化模型
    \item 相比基线方法，提升多视角一致性与真实感效果
    \item 负责算法实现、模型优化、实验对比
\end{itemize}

\section{\makebox[\widthof{\faFileTextO}][c]{\color{CVBlue}\faFileTextO}\ 论文成果}
\begin{itemize}[nosep]
    \item \textbf{LIFTINGAVATAR: HIGH-FIDELITY ANIMATABLE 3D GAUSSIAN AVATAR FROM SINGLE IMAGE} \\
    2026 | 第一作者 | 在投
\end{itemize}

\section{\makebox[\widthof{\faCogs}][c]{\color{CVBlue}\faCogs}\ 专业技能}
\begin{itemize}[nosep]
    \item \textbf{编程语言：} Python / C / MATLAB
    \item \textbf{框架工具：} PyTorch / Git / Linux / Arduino
    \item \textbf{研究方向：} 计算机视觉、3D重建、GNN、深度学习、嵌入式基础
\end{itemize}

\section{\makebox[\widthof{\faLanguage}][c]{\color{CVBlue}\faLanguage}\ 语言能力}
\begin{itemize}[nosep]
    \item 英语：IELTS 7.0
    \item 具备剑桥大学、UBC海外科研经历
\end{itemize}

\section{\makebox[\widthof{\faStar}][c]{\color{CVBlue}\faStar}\ 个人优势}
\begin{itemize}[nosep]
    \item 专业成绩优异，数学与算法基础扎实
    \item 拥有完整AI项目实践，独立完成模型设计、训练与实验验证
\end{itemize}

\end{document}
